\begin{problem}[Warmup]{(a)}
    If $[K : F] = 1$, what can you say about the fields $K$ and $F$?
\end{problem}
\unsolved

\begin{problem}[Warmup]{(b)}
    If $[K : F]$ is prime and $E$ is a field where $F \subseteq E \subseteq K$, what can you say about the field $E$?
\end{problem}
\unsolved

\begin{problem}[Warmup]{(c)}
    If $F \subseteq E \subseteq K$ is a sequence of finite field extensions and $[K : F] = [E : F]$, what can you say about the fields $K$ and $E$?
\end{problem}
\unsolved

\begin{problem}[Warmup]{(d)}
    Explain why the following statements are true, or else give a counterexample.
    \begin{enumerate}
        \item[(a)] Every finite extension is algebraic.
        \item[(b)] Every simple extension is algebraic.
        \item[(c)] Every simple extension is finite.
        \item[(d)] Every algebraic extension is simple.
        \item[(e)] Every algebraic extension is finite.
    \end{enumerate}
\end{problem}
\unsolved

\begin{problem}[Warmup]{(e)}
    How many elements belong to a field that is a degree 2 extension of a degree 3 extension of $\mathbb{Z}_7$?
\end{problem}
\unsolved

\begin{problem}[Warmup]{(f)}
    Give examples of the following, or else explain why the example does not exist:
    \begin{enumerate}
        \item[(a)] A degree 2 extension of $\mathbb{Q}(\sqrt{2})$.
        \item[(b)] A degree 3 extension of $\mathbb{R}$.
        \item[(c)] A degree 3 extension of $\mathbb{Z}_2$.
    \end{enumerate}
\end{problem}
\unsolved

\begin{problem}{1}
    Show that $\sqrt{2} \notin \mathbb{Q}(\sqrt{3})$, using Theorem 43.2. (See
    Example 38.2 for a more elementary solution of this exercise.)
\end{problem}
\unsolved

\begin{problem}{2}
    Prove Corollary 44.4. That is, suppose that each $F_i$ is a field, and
    $F_1 \subseteq F_2 \subseteq \cdots \subseteq F_n$ is a sequence of finite extensions.
    Show that
    \[[F_n : F_1] = [F_n : F_{n - 1}][F_{n - 1} : F_{n - 2}] \cdots [F_2 : F_1].\]
\end{problem}
\unsolved

\begin{problem}{3}
    Prove that $\mathbb{Q}(\sqrt{2} + \sqrt{3}) = \mathbb{Q}(\sqrt{2}, \sqrt{3})$.
\end{problem}
\unsolved

\begin{problem}{4}
    \begin{enumerate}
        \item[(a)] Determine $[\mathbb{Q}(\sqrt[3]{2}, i) : \mathbb{Q}]$; include a proof
        of your result.
        \item[(b)] Find $\alpha \in \mathbb{C}$ so that
        $\mathbb{Q}(\alpha) = \mathbb{Q}(\sqrt[3]{2}, i)$.
    \end{enumerate}
\end{problem}
\unsolved

\begin{problem}{5}
    In this exercise you will show that not every algebraic extension is a finite
    extension by considering the field $F$, constructed as follows. Let
    \[F_1 = \mathbb{Q}(\sqrt{2}),\ F_2 = F_1(\sqrt[3]{2}),\ F_3 = F_2(\sqrt[4]{2}),\
    F_4 = F_3(\sqrt[5]{2}),\ \ldots,\]
    and continue inductively. Then let
    \[F = \bigcup_{n = 1}^{\infty} F_n.\]
    Argue that $F$ is field. Then prove that $F$ is not a finite extension, thus showing
    that the converse of Theorem 44.1 is false.
\end{problem}
\unsolved

\begin{problem}{6}
    Give an example to show that if $[F(a) : F] = n$ and $[F(b) : F] = m$, it does not
    necessarily follow that $[F(a, b) : F] = mn$.
\end{problem}
\unsolved

\begin{problem}{7}
    Prove that an algebraic extension of an algebraic extension is an algebraic extension.
\end{problem}
\unsolved

\begin{problem}{8}
    Consider the primitive cube root of unity
    \[\zeta = -\frac{1}{2} + i\frac{\sqrt{3}}{2} = e^{\frac{2\pi i}{3}}
    = \cos\left(\frac{2\pi}{3}\right) + i\sin\left(\frac{2\pi}{3}\right).\]
    Compute $[\mathbb{Q}(\zeta) : \mathbb{Q}]$.
\end{problem}
\unsolved

\begin{problem}{9}
    Consider the field $F = \mathbb{Z}_2(\alpha, \beta)$ constructed in Exercise 43.2.
    \begin{enumerate}
        \item[(a)] Use Theorem 44.2 to determine the possible degrees $[F : K]$, where $K$
        is a subfield of $F$. Find a subfield with each of these degrees, and draw a field
        extension diagram similar to that in Examples 44.2 and 44.3.
        \item[(b)] Prove that $\mathbb{Z}_2(\alpha + \beta) = \mathbb{Z}_2(\alpha, \beta)$.
    \end{enumerate}
\end{problem}
\unsolved

\begin{problem}{10}
    Prove that $\mathbb{A}$, the set of all complex numbers algebraic over $\mathbb{Q}$, is
    a field.
\end{problem}
\unsolved
